\documentclass[a4paper,12pt]{article}

\usepackage[T2A]{fontenc}
\usepackage[utf8]{inputenc}
\usepackage[russian]{babel}
\usepackage{amsmath,amssymb,mathtools}
\usepackage{PTSans}
\renewcommand{\familydefault}{\sfdefault}
\usepackage{icomma}
\usepackage[left=20mm,right=20mm,top=20mm,bottom=22mm]{geometry}
\usepackage{microtype}
\usepackage{xcolor}
\usepackage{tikz}
\usetikzlibrary{calc,arrows.meta,positioning,shapes.geometric}
\usepackage{pgfplots}
\pgfplotsset{compat=1.18}
\usepackage{tabularx,booktabs}
\usepackage{enumitem}
\usepackage{parskip}
\usepackage{fancyhdr}
\usepackage{setspace}
\onehalfspacing

\definecolor{accent}{HTML}{008080}
\definecolor{darkaccent}{HTML}{004D4D}
\definecolor{textgray}{HTML}{333333}
\definecolor{softgray}{HTML}{6B7474}

\color{textgray}

\newcommand{\checkbox}{\(\square\)}

\newcommand{\topicrule}{%
  \par\noindent
  {\color{accent}\rule{\linewidth}{0.55pt}}
  \par
}

\newcommand{\keyformula}[1]{%
  \begin{center}
    {\color{darkaccent}\large #1}
  \end{center}
}

\setlist[itemize]{
  leftmargin=7mm,
  itemsep=3pt,
  topsep=3pt
}

\setlist[enumerate]{
  leftmargin=8mm,
  itemsep=4pt,
  topsep=4pt
}

\pagestyle{fancy}
\fancyhf{}
\lhead{\small\color{softgray}Николь Саркисьянц}
\rhead{\small\color{softgray}Производная и касательная}
\cfoot{\small\color{softgray}\thepage}
\renewcommand{\headrulewidth}{0.3pt}
\renewcommand{\footrulewidth}{0pt}
\setlength{\headheight}{14pt}

\tikzset{
  flowstart/.style={
    ellipse,
    draw=darkaccent,
    line width=0.8pt,
    minimum width=42mm,
    minimum height=11mm,
    align=center,
    font=\small
  },
  flowaction/.style={
    rectangle,
    rounded corners=2mm,
    draw=accent,
    line width=0.8pt,
    text width=100mm,
    minimum height=12mm,
    inner sep=3mm,
    align=center,
    font=\small
  },
  flowsmall/.style={
    rectangle,
    rounded corners=2mm,
    draw=accent,
    line width=0.8pt,
    text width=48mm,
    minimum height=12mm,
    inner sep=3mm,
    align=center,
    font=\small
  },
  flowdecision/.style={
    diamond,
    aspect=2.35,
    draw=darkaccent,
    line width=0.8pt,
    text width=42mm,
    inner sep=1.5mm,
    align=center,
    font=\small
  },
  flowarrow/.style={
    -{Latex[length=3mm,width=2mm]},
    line width=0.8pt,
    draw=darkaccent
  },
  flowlabel/.style={
    font=\small,
    fill=white,
    inner sep=1pt
  }
}

\begin{document}

% =========================================================
% Титульный лист
% =========================================================

\begin{titlepage}
\thispagestyle{empty}

\begin{center}

\vspace*{20mm}

{\Large\color{darkaccent}
Лёвин Артём Александрович
}

\vspace{20mm}

{\color{accent}\rule{0.70\linewidth}{0.8pt}}

\vspace{13mm}

{\Huge\bfseries
Чек-лист
}

\vspace{7mm}

{\LARGE\bfseries
Производная и касательная
}

\vspace{5mm}

{\Large
Угловой коэффициент. Графический смысл производной.\\
Задачи на касательную
}

\vspace{16mm}

{\large
ЕГЭ по профильной математике
}

\vspace{20mm}

{\Large\bfseries
Николь Саркисьянц
}

\vfill

{\small
Лёвин Артём Александрович\\
эксклюзивно для Николь Саркисьянц
}

\vspace{5mm}

{\small \today}

\end{center}

\begin{tikzpicture}[remember picture,overlay]
  \draw[
    accent!55,
    line width=1.1pt
  ]
  ([xshift=18mm,yshift=18mm]current page.south west)
  .. controls
  ([xshift=58mm,yshift=31mm]current page.south west)
  and
  ([xshift=-65mm,yshift=8mm]current page.south east)
  ..
  ([xshift=-18mm,yshift=20mm]current page.south east);
\end{tikzpicture}

\end{titlepage}

\setcounter{page}{1}

% =========================================================
% 1. Ядро темы
% =========================================================

\section*{\color{darkaccent}1. Ядро темы}
\topicrule

В задачах этого типа одно и то же число можно увидеть тремя способами:
как значение производной, как угловой коэффициент касательной и как
тангенс угла её наклона.

\keyformula{
\[
\boxed{f'(x_0)=k=\tg\alpha}
\]
}

Здесь:

\begin{itemize}
  \item[\checkbox]
  \(f'(x_0)\) --- производная функции в точке с абсциссой \(x_0\);

  \item[\checkbox]
  \(k\) --- угловой коэффициент касательной;

  \item[\checkbox]
  \(\alpha\) --- угол наклона касательной к положительному направлению
  оси \(Ox\), где
  \[
  0\le\alpha<\pi;
  \]

  \item[\checkbox]
  касательная имеет уравнение прямой
  \[
  y=kx+b.
  \]
\end{itemize}

Для стандартных задач ЕГЭ предполагается, что производная в рассматриваемой
точке существует и конечна. Тогда касательная не вертикальна и её угловой
коэффициент определён.

\section*{\color{darkaccent}2. Абсцисса и ордината}
\topicrule

\begin{itemize}
  \item[\checkbox]
  \textbf{Абсцисса} --- координата \(x\), то есть положение точки
  по горизонтали.

  \item[\checkbox]
  \textbf{Ордината} --- координата \(y\), то есть положение точки
  по вертикали.

  \item[\checkbox]
  Если дана точка с абсциссой \(x_0\), то в этой точке
  \[
  x=x_0.
  \]
\end{itemize}

\section*{\color{darkaccent}3. Как знак производной связан с касательной}
\topicrule

\begin{center}
\renewcommand{\arraystretch}{1.35}

\begin{tabularx}{0.94\linewidth}
{@{}>{\centering\arraybackslash}p{31mm}X@{}}

\toprule

\textbf{Знак}
&
\textbf{Что видно на касательной}
\\

\midrule

\(f'(x_0)>0\)
&
Касательная поднимается слева направо.
\\

\(f'(x_0)<0\)
&
Касательная опускается слева направо.
\\

\(f'(x_0)=0\)
&
Касательная горизонтальна.
\\

\bottomrule

\end{tabularx}
\end{center}

Самая надёжная вычислительная формула для прямой через две точки

\[
A(x_1;y_1),
\qquad
B(x_2;y_2)
\]

имеет вид

\[
\boxed{
k=
\frac{y_2-y_1}{x_2-x_1}
}
\qquad
(x_1\ne x_2).
\]

Эта формула автоматически учитывает знак наклона.

Если Вы строите прямоугольный треугольник и используете только
\emph{длины} его катетов, то сначала получается величина \(|k|\).
После этого знак определяется по направлению касательной.

% =========================================================
% 4. Производная по изображённой касательной
% =========================================================

\section*{\color{darkaccent}4. Производная по изображённой касательной}
\topicrule

Чтобы найти \(f'(x_0)\) по рисунку:

\begin{enumerate}
  \item
  найдите касательную, проведённую в точке с абсциссой \(x_0\);

  \item
  выберите на касательной две точки с точно читаемыми координатами;

  \item
  вычислите её угловой коэффициент
  \[
  k=
  \frac{\Delta y}{\Delta x};
  \]

  \item
  запишите
  \[
  f'(x_0)=k.
  \]
\end{enumerate}

\textbf{Важно.}
Выбранные точки должны лежать именно на касательной.
Точка касания при этом не обязана быть одной из выбранных точек.

\subsection*{Простой пример}

Рассмотрим функцию

\[
f(x)=\frac{x^2}{4}.
\]

В точке

\[
x_0=2
\]

имеем

\[
f(2)=1,
\]

а касательная имеет уравнение

\[
y=x-1.
\]

На касательной удобно взять точки

\[
(0;-1)
\qquad\text{и}\qquad
(2;1).
\]

Тогда

\[
k=
\frac{1-(-1)}{2-0}
=
\frac22
=
1.
\]

Следовательно,

\[
\boxed{f'(2)=1}.
\]

\begin{center}

\begin{tikzpicture}

\begin{axis}[
  width=100mm,
  height=100mm,
  axis lines=middle,
  xmin=-2.5,
  xmax=4.5,
  ymin=-2,
  ymax=5,
  xlabel={$x$},
  ylabel={$y$},
  xtick={-2,-1,0,1,2,3,4},
  ytick={-1,0,1,2,3,4},
  samples=220,
  clip=false,
  unit vector ratio=1 1,
  grid=both,
  grid style={
    line width=.15pt,
    draw=gray!25
  },
  axis line style={
    draw=textgray
  },
  tick label style={
    font=\small
  }
]

\addplot[
  domain=-2.3:4.3,
  line width=1.2pt,
  color=darkaccent
]
{x^2/4};

\addplot[
  domain=-0.8:4.3,
  line width=1.1pt,
  color=accent
]
{x-1};

\addplot[
  only marks,
  mark=*,
  mark size=2pt,
  color=darkaccent
]
coordinates {(2,1)};

\addplot[
  only marks,
  mark=*,
  mark size=1.8pt,
  color=accent
]
coordinates {(0,-1) (2,1)};

\draw[
  dashed,
  accent
]
(axis cs:0,-1)
--
(axis cs:2,-1);

\draw[
  dashed,
  accent
]
(axis cs:2,-1)
--
(axis cs:2,1);

\node[
  anchor=south west,
  font=\small,
  text=darkaccent
]
at (axis cs:2,1)
{$(2;1)$};

\node[
  anchor=north east,
  font=\small,
  text=accent
]
at (axis cs:0,-1)
{$(0;-1)$};

\node[
  anchor=south west,
  font=\small,
  text=darkaccent
]
at (axis cs:3.0,2.5)
{$y=\dfrac{x^2}{4}$};

\node[
  anchor=south west,
  font=\small,
  text=accent
]
at (axis cs:3.0,1.75)
{$y=x-1$};

\end{axis}

\end{tikzpicture}

\end{center}

\subsection*{Если касательная убывает}

Пусть касательная проходит через точки

\[
A(-1;3),
\qquad
B(3;1).
\]

Тогда

\[
k=
\frac{1-3}{3-(-1)}
=
-\frac24
=
-\frac12.
\]

Поэтому

\[
\boxed{
f'(x_0)=-\frac12
}.
\]

% =========================================================
% 5. Параллельность
% =========================================================

\section*{\color{darkaccent}5. Если касательная параллельна заданной прямой}
\topicrule

Параллельные прямые вида

\[
y=kx+b
\]

имеют одинаковые угловые коэффициенты.

Поэтому если касательная к графику функции в точке \(x_0\)
параллельна прямой

\[
y=mx+c,
\]

то

\[
\boxed{
f'(x_0)=m
}.
\]

Свободный коэффициент \(c\) на значение производной здесь не влияет.

\subsection*{Пример}

Касательная параллельна прямой

\[
y=4x-7.
\]

Её угловой коэффициент равен \(4\), следовательно,

\[
\boxed{
f'(x_0)=4
}.
\]

% =========================================================
% 6. График производной
% =========================================================

\section*{\color{darkaccent}6. Если дан график производной \(y=f'(x)\)}
\topicrule

На графике производной значение \(f'(x)\) является обычной ординатой \(y\).

Поэтому условие

\[
f'(x)=m
\]

на таком рисунке означает

\[
y=m.
\]

Отсюда:

\begin{itemize}
  \item[\checkbox]
  если нужно найти \textbf{количество} подходящих значений \(x\),
  посчитайте точки пересечения графика \(y=f'(x)\) с горизонталью
  \[
  y=m;
  \]

  \item[\checkbox]
  если нужны \textbf{абсциссы}, считайте координаты \(x\)
  соответствующих точек пересечения.
\end{itemize}

Типовая цепочка рассуждений:

\[
\boxed{
\text{параллельность}
\Rightarrow
k=m
\Rightarrow
f'(x)=m
\Rightarrow
y=m
\text{ на графике }f'
}
\]

% =========================================================
% 7. Касательная через заданную точку
% =========================================================

\section*{\color{darkaccent}7. Если касательная проходит через заданную точку}
\topicrule

Пусть известна точка касания

\[
M(x_0;f(x_0))
\]

и ещё одна точка касательной

\[
P(x_1;y_1),
\qquad
x_1\ne x_0.
\]

Тогда

\[
\boxed{
f'(x_0)
=
k
=
\frac{f(x_0)-y_1}{x_0-x_1}
}.
\]

Если касательная проходит через начало координат

\[
O(0;0)
\]

и

\[
x_0\ne0,
\]

то

\[
\boxed{
f'(x_0)
=
\frac{f(x_0)}{x_0}
}.
\]

% =========================================================
% 8. Функция и конкретная касательная
% =========================================================

\section*{\color{darkaccent}8. Если даны функция и конкретная касательная}
\topicrule

Пусть прямая

\[
y=kx+b
\]

является касательной к графику

\[
y=f(x)
\]

в точке с абсциссой \(x_0\).

В точке касания одновременно выполняются два условия:

\[
\boxed{
f'(x_0)=k
}
\]

и

\[
\boxed{
f(x_0)=kx_0+b
}.
\]

Первое условие означает совпадение наклонов.

Второе условие означает, что точка касания принадлежит одновременно
графику функции и касательной.

\subsection*{Пример из ключевого типа задач}

Прямая

\[
y=-2x-12
\]

является касательной к графику функции

\[
f(x)=x^3-2x^2-6x-4.
\]

Найдём абсциссу точки касания.

\textbf{1. Совпадение наклонов}

\[
f'(x)=3x^2-4x-6.
\]

Угловой коэффициент касательной равен

\[
k=-2.
\]

Следовательно,

\[
3x^2-4x-6=-2,
\]

откуда

\[
3x^2-4x-4=0.
\]

Дискриминант:

\[
D=
(-4)^2
-
4\cdot3\cdot(-4)
=
16+48
=
64.
\]

Получаем два кандидата:

\[
x_1=2,
\qquad
x_2=-\frac23.
\]

\textbf{2. Проверка принадлежности касательной}

Для \(x=2\):

\[
f(2)
=
8-8-12-4
=
-16,
\]

а на прямой

\[
-2\cdot2-12
=
-16.
\]

Кандидат подходит.

Для

\[
x=-\frac23
\]

получаем

\[
f\!\left(-\frac23\right)
=
-\frac{8}{27}
-
\frac89
+
4
-
4
=
-\frac{32}{27}.
\]

На касательной:

\[
-2\left(-\frac23\right)-12
=
\frac43-12
=
-\frac{32}{3}.
\]

Значения различаются, поэтому

\[
x=-\frac23
\]

не является абсциссой точки касания с данной прямой.

Итак,

\[
\boxed{
x_0=2
}.
\]

% =========================================================
% Блок-схема алгоритма
% =========================================================

\clearpage
\thispagestyle{empty}

\begin{center}

{\Large\bfseries\color{darkaccent}
Алгоритм: функция и уравнение касательной
}

\end{center}

\vspace{7mm}

\begin{center}

\begin{tikzpicture}[
  node distance=7mm and 22mm
]

\node[
  flowstart
]
(start)
{Даны \(y=f(x)\) и \(y=kx+b\)};

\node[
  flowaction,
  below=of start
]
(deriv)
{Найдите производную \(f'(x)\).};

\node[
  flowaction,
  below=of deriv
]
(slope)
{Выделите угловой коэффициент касательной \(k\).};

\node[
  flowaction,
  below=of slope
]
(candidates)
{
Решите уравнение
\[
f'(x)=k
\]
и получите кандидатов на абсциссу точки касания.
};

\node[
  flowaction,
  below=of candidates
]
(take)
{Возьмите очередного кандидата \(x\).};

\node[
  flowdecision,
  below=9mm of take
]
(check)
{
Выполняется ли
\(f(x)=kx+b\)?
};

\node[
  flowsmall,
  below right=11mm and 10mm of check
]
(keep)
{Сохраните этот кандидат.};

\node[
  flowsmall,
  below left=11mm and 10mm of check
]
(reject)
{Отбросьте этот кандидат.};

\node[
  flowdecision,
  below=24mm of check
]
(more)
{Остались непроверенные кандидаты?};

\node[
  flowstart,
  below=12mm of more
]
(finish)
{Оставшиеся значения \(x\) --- ответ};

\draw[flowarrow]
(start)
--
(deriv);

\draw[flowarrow]
(deriv)
--
(slope);

\draw[flowarrow]
(slope)
--
(candidates);

\draw[flowarrow]
(candidates)
--
(take);

\draw[flowarrow]
(take)
--
(check);

\draw[flowarrow]
(check.east)
-|
node[
  flowlabel,
  pos=0.25
]
{да}
(keep.north);

\draw[flowarrow]
(check.west)
-|
node[
  flowlabel,
  pos=0.25
]
{нет}
(reject.north);

\draw[flowarrow]
(keep.south)
|-
(more.east);

\draw[flowarrow]
(reject.south)
|-
(more.west);

\draw[flowarrow]
(more.south)
--
node[
  flowlabel,
  right
]
{нет}
(finish.north);

\draw[flowarrow]
(more.west)
--
++(-24mm,0)
|-
node[
  flowlabel,
  pos=0.25
]
{да}
(take.west);

\end{tikzpicture}

\end{center}

\clearpage

% =========================================================
% 9. Что знать наизусть
% =========================================================

\section*{\color{darkaccent}9. Что необходимо знать наизусть}
\topicrule

\begin{itemize}

  \item[\checkbox]
  Касательная в рассматриваемой точке имеет вид
  \[
  y=kx+b.
  \]

  \item[\checkbox]
  Основная связь:
  \[
  f'(x_0)=k=\tg\alpha.
  \]

  \item[\checkbox]
  По двум точкам прямой:
  \[
  k=
  \frac{y_2-y_1}{x_2-x_1}.
  \]

  \item[\checkbox]
  Параллельные прямые имеют одинаковые угловые коэффициенты.

  \item[\checkbox]
  Если касательная параллельна
  \[
  y=mx+c,
  \]
  то
  \[
  f'(x_0)=m.
  \]

  \item[\checkbox]
  На графике
  \[
  y=f'(x)
  \]
  условие
  \[
  f'(x)=m
  \]
  означает
  \[
  y=m.
  \]

  \item[\checkbox]
  Для конкретной касательной
  \[
  y=kx+b
  \]
  в точке касания одновременно выполняются условия
  \[
  f'(x_0)=k,
  \qquad
  f(x_0)=kx_0+b.
  \]

\end{itemize}

% =========================================================
% 10. Типичные ошибки
% =========================================================

\section*{\color{darkaccent}10. Типичные ошибки}
\topicrule

\begin{enumerate}

  \item
  \textbf{Брать две точки на графике функции вместо касательной.}

  Для вычисления \(k\) нужны две точки одной и той же касательной.

  \item
  \textbf{Использовать неточные точки.}

  Лучше выбирать узлы координатной сетки, координаты которых считываются
  однозначно.

  \item
  \textbf{Терять знак производной.}

  Если касательная опускается слева направо, то

  \[
  k<0.
  \]

  \item
  \textbf{Перепутать вертикальное и горизонтальное изменения.}

  Всегда

  \[
  k=
  \frac{\Delta y}{\Delta x}.
  \]

  \item
  \textbf{Брать свободный коэффициент \(b\) вместо \(k\).}

  С производной связан коэффициент перед \(x\).

  \item
  \textbf{Перепутать график функции и график производной.}

  На графике

  \[
  y=f'(x)
  \]

  значение производной читается как координата \(y\).

  \item
  \textbf{После решения \(f'(x)=k\) принимать все корни без проверки.}

  Если дана конкретная касательная

  \[
  y=kx+b,
  \]

  каждый кандидат должен удовлетворять также условию

  \[
  f(x)=kx+b.
  \]

  \item
  \textbf{Считать, что после нахождения одного подходящего кандидата
  остальные можно не проверять.}

  В общем случае одна прямая может касаться графика более чем в одной точке,
  поэтому следует проверить все найденные кандидаты.

\end{enumerate}

% =========================================================
% 11. Краткая памятка
% =========================================================

\section*{\color{darkaccent}11. Сверхкраткая памятка}
\topicrule

\begin{center}

\renewcommand{\arraystretch}{1.5}

\begin{tabularx}{0.95\linewidth}
{@{}p{54mm}X@{}}

\toprule

\textbf{Что дано}
&
\textbf{Что делать}
\\

\midrule

Нарисована касательная
&
Взять две точные точки на ней и найти
\[
k=
\frac{\Delta y}{\Delta x}.
\]
\\

Касательная параллельна
\(y=mx+c\)
&
Сразу использовать
\[
f'(x_0)=m.
\]
\\

Дан график
\(y=f'(x)\)
и известно
\(f'(x)=m\)
&
Искать на графике точки с ординатой
\[
y=m.
\]
\\

Даны
\(f(x)\)
и касательная
\(y=kx+b\)
&
Сначала решить
\[
f'(x)=k,
\]
затем проверить кандидатов по
\[
f(x)=kx+b.
\]
\\

Касательная убывает
&
Помнить:
\[
f'(x_0)<0.
\]
\\

Касательная горизонтальна
&
\[
f'(x_0)=0.
\]
\\

\bottomrule

\end{tabularx}

\end{center}



\end{document}